\chapter{Detalle del protocolo y verificaciones}
\label{anx:protocolo}
\small

\section{Cálculo de la suma de comprobación}

El alfabeto Contact ID asigna a los dígitos 1--9 su valor, al 0 el valor 10 y a B--F los valores
11--15. La suma de comprobación $S$ satisface $(\sum v(d_i) + v(S)) \bmod 15 = 0$. Ejemplo verificado:
cuenta 1234, nuevo evento 381 (pérdida de supervisión RF), zona 007: la suma del cuerpo es 79;
$79 \bmod 15 = 4$; falta 11, cuyo símbolo es B. Mensaje: \texttt{123418138100007B}.

\section{Ejemplos generados y verificados programáticamente}

Producidos con la herramienta propia \path{plataforma/herramientas/contactid.py}, que genera,
valida y empaqueta mensajes; la alteración de un solo dígito invalida la suma de comprobación.

\begin{longtable}{@{}L{4.4cm}L{4.2cm}L{3.6cm}@{}}
\filaencab \encab{Evento} & \encab{Mensaje} & \encab{BCD (8 bytes)} \\ \midrule \endhead
Nodo 007 dejó de reportar & \texttt{1234 18 1 381 00 007 B} & \texttt{12 34 18 13 81 00 00 7B} \\
Nodo 007 volvió a reportar & \texttt{1234 18 3 381 00 007 9} & \texttt{12 34 18 33 81 00 00 79} \\
Batería baja en nodo 012 & \texttt{1234 18 1 384 00 012 7} & \texttt{12 34 18 13 84 00 01 27} \\
Temperatura alta, zona 003 & \texttt{1234 18 1 158 01 003 7} & \texttt{12 34 18 11 58 01 00 37} \\
Reporte periódico & \texttt{1234 18 1 602 00 000 2} & \texttt{12 34 18 16 02 00 00 02} \\
Interferencia RF en gateway & \texttt{1234 18 1 344 00 000 9} & \texttt{12 34 18 13 44 00 00 09} \\
\bottomrule
\caption{Mensajes de ejemplo con suma de comprobación válida.}
\end{longtable}

\section{Tamaño del mensaje unificado}

\begin{longtable}{@{}L{6.0cm}C{1.9cm}C{1.9cm}C{2.0cm}@{}}
\filaencab \encab{Mensaje} & \encab{Evento} & \encab{Valor} & \encab{Total} \\ \midrule \endhead
Solo evento & 8 B & --- & \textbf{8 B} \\
Evento con batería \texttt{[B12]} & 8 B & 5 B & 13 B \\
Evento con temperatura \texttt{[T23.4]} & 8 B & 7 B & 15 B \\
Evento con RSSI/SNR \texttt{[S-103/4.2]} & 8 B & 11 B & 19 B \\
\bottomrule
\caption{Verificación programática del tamaño. El payload LoRaWAN disponible es de 11 bytes en la tasa
más lenta y de 51--53 bytes en operación normal: por eso el valor es opcional.}
\end{longtable}

\section{Identificadores del campo de datos extendidos (perfil CID-IoT)}

\begin{longtable}{@{}C{1.4cm}L{7.4cm}L{3.4cm}@{}}
\filaencab \encab{Id.} & \encab{Contenido} & \encab{Ejemplo} \\ \midrule \endhead
\texttt{T} & Temperatura en grados Celsius & \texttt{[T23.4]} \\
\texttt{H} & Humedad relativa en por ciento & \texttt{[H67]} \\
\texttt{N} & Nivel en por ciento o centímetros & \texttt{[N45]} \\
\texttt{B} & Batería en por ciento & \texttt{[B87]} \\
\texttt{V} & Tensión en voltios & \texttt{[V3.28]} \\
\texttt{S} & Calidad de señal RSSI/SNR & \texttt{[S-103/4.2]} \\
\texttt{Q} & Contador de secuencia & \texttt{[Q1042]} \\
\bottomrule
\caption{Identificadores reservados. La tabla se cierra en la Etapa 2.}
\end{longtable}

\section{Mapeo de sensores heterogéneos a eventos normalizados (extracto)}

\begin{longtable}{@{}L{2.0cm}L{4.2cm}L{3.0cm}C{1.2cm}C{0.8cm}@{}}
\filaencab \encab{Transporte} & \encab{Señal} & \encab{Condición} & \encab{Código} & \encab{Q} \\ \midrule \endhead
LoRaWAN & Sin transmisión en 3 intervalos & Pérdida de supervisión & 381 & 1 \\
LoRaWAN & Batería bajo umbral & Batería baja & 384 & 1 \\
MQTT & \emph{Last Will} recibido & Cliente desconectado & 381 & 1 \\
MQTT/HTTP & Temperatura sobre umbral & Temperatura alta & 158 & 1 \\
WiFi/GPIO & Contacto magnético abierto & Apertura de acceso & 134 & 1 \\
Panel de alarma & Ya emite Contact ID & Sin traducción & --- & --- \\
\bottomrule
\caption{Extracto del mapeo del perfil. La tabla completa se desarrolla en la Etapa 2; la versión de
trabajo está en el repositorio (\texttt{docs/05-protocolo-contactid.md}).}
\end{longtable}
